% Part: turing-machines
% Chapter: undecidability
% Section: representing-tms

\documentclass[../../../include/open-logic-section]{subfiles}

\begin{document}

\olfileid{tur}{und}{rep}
\olsection{د ټيورينګ ماشينونو تمثيلول}

\begin{explain}
د لومړۍ درجې منطق په يوه !!a{sentence} د ټيورينګ ماشينونو او د هغو
د چلند د تمثيلولو دپاره بايد يوه مناسبه ژبه تعريف کړو۔ ژبه دوه برخې
لري: د ماشين د تشکيلونو د بيانولو دپاره !!{predicate}s، او د اجرا د
ګامونو (``شېبو'') او پر پټه د موقعيتونو د شمېرلو دپاره تعبيرونه۔

موږ دوه ډوله !!{predicate}s ورپېژنو چې دواړه دوه-ځايه دي: د هر
حالت~$q$ دپاره يو !!a{predicate}~$\Obj Q_q$، او د پټې د هرې
نښې~$\sigma$ دپاره يو !!a{predicate}~$\Obj S_\sigma$۔ لومړي موږ ته
اجازه راکوي چې د~$M$ حالت او د هغه د پټې د سر موقعيت بيان کړو، او
دويم موږ ته اجازه راکوي چې د پټې منځپانګه بيان کړو۔

د پټې د سر د موقعيتونو او د اجرا شوو ګامونو د شمېر د څرګندولو دپاره
د اعدادو د څرګندولو يوې لارې ته اړ يو۔ دا د يوه !!a{constant}~$\Obj
0$، او د يوې $1$-ځايه تابع~$\prime$، يعنې د تالي تابع، په وسيله کېږي۔
د دود له مخې دا د خپل آرګومان \emph{وروسته} ليکل کېږي (او قوسونه
ترې پرېږدو)۔

د هر عدد~$n$ دپاره يو معياري ترم~$\num{n}$ شته، يعنې د~$n$
\emph{عددنښه}، چې په~$\Lang L_M$ کښې ئې تمثيلوي۔ $\num{0}$ هم
$\Obj 0$ دے، $\num{1}$ هم~$\Obj 0'$ دے، $\num{2}$ هم~$\Obj 0''$
دے، او داسې نور۔ په لا صوري ډول:
\begin{align*}
\num{0} & = \Obj 0 \\
\num{n+1} &= \num{n}'
\end{align*}
ترم~$\num{0}$، يعنې~$\Obj 0$، هم پر پټه تر ټولو کيڼ موقعيت او هم د
اجرا له لومړي ګام څخه مخکې وخت (پيلنی تشکيل) نوموي۔ ترم~$\num{1}$،
يعنې~$\Obj 0'$، د تر ټولو کيڼې مربع ښي لور ته مربع او د اجرا له لومړي
ګام څخه وروسته وخت نوموي، او داسې نور۔

موږ يو !!a{predicate}~$<$ هم ورپېژنو، څو هم د پټې د موقعيتونو ترتيب
(چې هلته معنا ئې ``کيڼ لور ته'' ده) او هم د اجرا د ګامونو ترتيب
(چې هلته معنا ئې ``مخکې'' ده) څرګند کړو۔

کله چې ژبه مو تياره شوه، د~$!T(M, w)$ ``اکسيومونه'' لست کوو؛ يعنې
هغه !!{sentence}s چې په ګډه د~$M$ چلند بيانوي کله چې پر ننوت~$w$
وچلول شي۔ داسې !!{sentence}s به وي چې پر~$\Obj 0$، $\prime$، او~$<$
شرطونه ږدي، داسې !!{sentence}s چې د ننوت تشکيل بيانوي، او داسې
!!{sentence}s چې د~$M$ تشکيل د يوې ځانګړې لارښوونې له اجرا وروسته
بيانوي۔
\end{explain}

\begin{defn}
  \ollabel{defn:tm-descr}
يو ټيورينګ ماشين~$M = \tuple{Q, \Sigma, q_0, \delta}$ چې راکړل شي،
ژبه~$\Lang L_M$ له دې شيانو جوړه ده:
\begin{enumerate}
\item د هر حالت~$q \in Q$ دپاره يو دوه-ځايه !!{predicate}
  $\Obj Q_q(x, y)$۔ په شهودي ډول~$\Obj Q_q(\num{m}, \num{n})$ دا
  څرګندوي: ``له~$n$ ګامونو وروسته، $M$ په حالت~$q$ کښې دے او~$m$-مه
  مربع لولي۔''
\item د هرې نښې~$\sigma\in \Sigma$ دپاره يو دوه-ځايه !!{predicate}
  $\Obj S_\sigma(x, y)$۔ په شهودي ډول~$\Obj S_\sigma(\num{m},
  \num{n})$ دا څرګندوي: ``له~$n$ ګامونو وروسته، $m$-مه مربع نښه
  $\sigma$ لري۔''
\item يو !!{constant}~$\Obj 0$
\item يوه يو-ځايه !!{function}~$\prime$
\item يو دوه-ځايه !!{predicate}~$<$
\end{enumerate}
\end{defn}

هغه !!{sentence}s چې پر ننوت~$w = \sigma_{i_1}\dots\sigma_{i_k}$
د ټيورينګ ماشين~$M$ عمل بيانوي دا دي:
\begin{enumerate}
\item هغه اکسيومونه چې اعداد او~$<$ بيانوي:
\begin{enumerate}
\item يو !!^a{sentence} چې وايي هر عدد له خپل تالي څخه کوچنی دے:
\[
\lforall[x][x < x']
\]
\item يو !!^a{sentence} چې د~$<$ تعدي تضمينوي:
\[
\lforall[x][\lforall[y][\lforall[z][
      ((x < y \land y < z) \lif x < z)]]]
\]
\end{enumerate}
\item هغه اکسيومونه چې د ننوت تشکيل بيانوي:
\begin{enumerate}
\item له~$0$ ګامونو وروسته---مخکې له دې چې ماشين پيل شي---$M$ په
  پيلني حالت~$q_0$ کښې دے او مربع~$1$ لولي:
\[
\Obj Q_{q_0}(\num{1}, \num{0})
\]
\item لومړۍ $k+1$ مربعې نښې~$\TMendtape$، $\sigma_{i_1}$، \dots،
  $\sigma_{i_k}$ لري:
\[
\Obj S_\TMendtape(\num{0}, \num{0}) \land
\Obj S_{\sigma_{i_1}}(\num{1}, \num{0}) \land
\dots \land
\Obj S_{\sigma_{i_k}}(\num{k}, \num{0})
\]
\item له دې پرته پټه تشه ده:
\[
\lforall[x][(\num{k} < x \lif \Obj S_\TMblank(x, \num{0}))]
\]
\end{enumerate}
\item هغه اکسيومونه چې له يوه تشکيل څخه بل ته د حالت بدلون بيانوي:

د لاندې برخې دپاره، $!A(x, y)$ د ټولو هغو !!{sentence}s عطف وي چې
دا بڼه لري:
\[
\lforall[z][
  (((z < x \lor x < z) \land \Obj S_\sigma(z, y))
  \lif \Obj S_\sigma(z, y'))]
\]
چې~$\sigma \in \Sigma$ وي۔ موږ~$!A(\num{m},\num{n})$ د دې څرګندولو
دپاره کاروو: ``له مربع~$m$ پرته، له~$n+1$ ګامونو وروسته پټه هماغسې
ده لکه له~$n$ ګامونو وروسته۔''
\begin{enumerate}
\item \ollabel{rep-right} د هرې لارښوونې~$\delta(q_i, \sigma) =
  \tuple{q_j, \sigma', \TMright}$ دپاره دا !!{sentence}:
\begin{align*}
& \lforall[x][\lforall[y][(
   (\Obj Q_{q_i}(x, y) \land \Obj S_{\sigma}(x, y)) \lif {}]] \\
&\qquad   (\Obj Q_{q_j}(x', y') \land \Obj S_{\sigma'}(x, y') \land
!A(x, y)))
\end{align*}
دا وايي چې که ماشين له~$y$ ګامونو وروسته په حالت~$q_i$ کښې وي او
هغه مربع~$x$ لولي چې نښه~$\sigma$ لري، نو له~$y+1$ ګامونو وروسته
مربع~$x+1$ لولي، په حالت~$q_j$ کښې وي، مربع~$x$ اوس~$\sigma'$ لري،
او له~$x$ پرته هره مربع هماغه نښه لري چې له~$y$ ګامونو وروسته ئې
درلوده۔

\item \ollabel{rep-left} د هرې لارښوونې~$\delta(q_i, \sigma) =
  \tuple{q_j, \sigma', \TMleft}$ دپاره دا !!{sentence}:
\begin{align*}
& \lforall[x][\lforall[y][
    ((\Obj Q_{q_i}(x', y) \land \Obj S_{\sigma}(x', y)) \lif {}]]\\
& \qquad   (\Obj Q_{q_j}(x, y') \land \Obj S_{\sigma'}(x', y') \land
!A(x', y))) \land {}\\
& \lforall[y][((\Obj Q_{q_i}(\num{0}, y) \land \Obj S_{\sigma}(\num{0},
    y)) \lif {}]\\
& \qquad (\Obj Q_{q_j}(\num{0}, y') \land \Obj S_{\sigma'}(\num{0},
  y') \land !A(\num{0}, y)))
\end{align*}
يوه شېبه فکر وکړئ چې دا څنګه کار کوي: اوس له ``که مربع~$x$ لولي
\dots'' څخه نۀ، بلکې له ``که مربع~$x+1$ لولي~\dots'' څخه پيل
کوو۔ کيڼ لور ته حرکت معنا دا چې په بل ګام کښې ماشين مربع~$x$ لولي۔
خو هغه مربع~$x+1$ ده چې څه پرې ليکل کېږي۔ دا کار ځکه داسې کوو چې
منفي کول يا د مخکيني تابع نۀ لرو۔
\textbf{د ژباړې د نسخې سمون (OLCMP-051):} د سرچينې په لومړي
کيڼ-حرکت اکسيوم کښې د نابدلې پټې شرط ناسمه مربع مستثنا کوله؛ دلته
هغه مربع مستثنا شوې چې ماشين پرې ليکي، څو له سملاسي ورپسې تشريح سره
او د حالت بدلون له معنا سره سمون ولري۔

په ياد ولرئ چې د~$x+1$ په بڼه اعداد~$1$، $2$، \dots، دي؛ يعنې دا
هغه حالت نۀ رانغاړي چې ماشين مربع~$0$ لولي او بايد کيڼ لور ته وخوځېږي
(چې البته نۀ شي کولے---يوازې پر ځاے پاتې کېږي)۔ هغه ځانګړے حالت د
دويم عطف له خوا رانغاړل کېږي: وايي چې که ماشين له~$y$ ګامونو وروسته
په حالت~$q_i$ کښې مربع~$0$ لولي او مربع~$0$ نښه~$\sigma$ لري، نو
له~$y+1$ ګامونو وروسته لا هم مربع~$0$ لولي، اوس په حالت~$q_j$ کښې
دے، پر مربع~$0$ نښه~$\sigma'$ ده، او له مربع~$0$ پرته مربعې هماغه
نښې لري چې له~$y$ ګامونو وروسته ئې لرلې۔
\textbf{د ژباړې د نسخې سمون (OLCMP-052):} په سرچينه کښې د ``له
ګامونو وروسته'' پر ځاے يوه چاپي تېروتنه راغلې وه؛ دلته سمه شوې ده۔
\item \ollabel{rep-stay} د هرې لارښوونې~$\delta(q_i, \sigma) =
  \tuple{q_j, \sigma', \TMstay}$ دپاره دا !!{sentence}:
\begin{align*}
& \lforall[x][\lforall[y][(
   (\Obj Q_{q_i}(x, y) \land \Obj S_{\sigma}(x, y)) \lif {}]] \\
&\qquad   (\Obj Q_{q_j}(x, y') \land \Obj S_{\sigma'}(x, y') \land
!A(x, y)))
\end{align*}
\end{enumerate}
\end{enumerate}
$!T(M, w)$ د ټيورينګ ماشين~$M$ او ننوت~$w$ دپاره د پورتنيو ټولو
!!{sentence}s عطف وي۔

د دې څرګندولو دپاره چې~$M$ بالاخره درېږي، بايد داسې !!{sentence}
ومومو چې وايي: ``له يو شمېر ګامونو وروسته به د حالت بدلون تابع
ناتعريفه وي۔'' $X$ د ټولو هغو جوړو~$\tuple{q, \sigma}$ سټ وي چې
$\delta(q, \sigma)$ ناتعريفه وي۔ بيا~$!E(M, w)$ دا !!{sentence} وي:
\[
\lexists[x][\lexists[y][(\bigvee_{\tuple{q, \sigma} \in
      X}(\Obj Q_q(x, y) \land \Obj S_\sigma(x, y)))]]
\]

که داسې ټيورينګ ماشين وکاروو چې ټاکل شوے درېدنی حالت~$h$ ولري، خبره
لا اسانه ده: بيا دا !!{sentence}~$!E(M, w)$
\[
\lexists[x][\lexists[y][\Obj Q_h(x, y)]]
\]
دا څرګندوي چې ماشين بالاخره درېږي۔

\begin{prop}
\ollabel{prop:mlessk}
که~$m < k$ وي، نو~$!T(M, w) \Entails \num{m} < \num{k}$۔
\end{prop}

\begin{proof}
تمرین۔
\end{proof}

\begin{prob}
\olref[tur][und][rep]{prop:mlessk} ثابت کړئ۔
(اشاره: پر~$k-m$ استقرا وکاروئ)۔
\end{prob}

\end{document}
